\documentclass[11pt]{article}

% ---------- packages ----------
\usepackage[T1]{fontenc}
\usepackage[utf8]{inputenc}
\usepackage{lmodern}
\usepackage{amssymb}
\usepackage[margin=1in]{geometry}
\usepackage{microtype}
\usepackage[dvipsnames,table]{xcolor}
\usepackage{titlesec}
\usepackage{enumitem}
\usepackage{booktabs}
\usepackage{array}
\usepackage{tabularx}
\usepackage{longtable}
\usepackage{tcolorbox}
\usepackage{fancyhdr}
\usepackage[hidelinks]{hyperref}

% ---------- colors ----------
\definecolor{primary}{HTML}{1F3A5F}   % deep blue
\definecolor{accent}{HTML}{8C5A2B}    % warm brown
\definecolor{light}{HTML}{EEF2F7}     % pale blue-gray
\definecolor{rule}{HTML}{C9D3DE}

\hypersetup{colorlinks=true, linkcolor=primary, urlcolor=accent}

% ---------- section styling ----------
\titleformat{\section}
  {\color{primary}\Large\bfseries}{\thesection.}{0.5em}{}
\titleformat{\subsection}
  {\color{accent}\large\bfseries}{\thesubsection}{0.5em}{}
\titlespacing*{\section}{0pt}{1.4em}{0.6em}

% pillar tag command
\newcommand{\pillar}[1]{{\small\color{accent}\textbf{[#1]}}}
\newcommand{\newtag}{{\small\color{accent}$\bigstar$\,new}}

% phase heading
\newcommand{\phase}[1]{%
  \vspace{0.8em}
  \noindent\colorbox{primary}{\parbox{\dimexpr\textwidth-2\fboxsep}{%
    \color{white}\bfseries\large\strut #1}}\par\vspace{0.5em}}

% week block
\newenvironment{week}[2]{%
  \vspace{0.4em}
  \noindent{\color{primary}\bfseries #1}\; #2\par\nobreak
  \begin{itemize}[leftmargin=1.4em,itemsep=1pt,topsep=2pt]}%
  {\end{itemize}}

% ---------- header/footer ----------
\pagestyle{fancy}
\fancyhf{}
\renewcommand{\headrulewidth}{0.4pt}
\renewcommand{\headrule}{\hbox to\headwidth{\color{rule}\leaders\hrule height \headrulewidth\hfill}}
\fancyhead[L]{\small\color{primary}AI for Theological Inquiry}
\fancyhead[R]{\small\color{primary}Mentored Research Seminar}
\fancyfoot[C]{\small\color{primary}\thepage}

\setlength{\parindent}{0pt}
\setlength{\parskip}{0.4em}

\renewcommand{\arraystretch}{1.25}
\newcolumntype{Y}{>{\raggedright\arraybackslash}X}

% ---------- title ----------
\begin{document}

\begin{center}
  {\color{primary}\Huge\bfseries AI for Theological Inquiry}\\[0.3em]
  {\color{accent}\Large Mentored Research Seminar}\\[0.9em]
  \begin{tcolorbox}[colback=light,colframe=rule,boxrule=0.5pt,arc=2pt,
    left=8pt,right=8pt,top=6pt,bottom=6pt,width=0.95\textwidth]
    \small
    \textbf{Format:} 14 weeks $\cdot$ one 100-minute session/week \quad
    \textbf{Type:} Mentored research seminar (students produce an original agentic-AI research project)\\[0.3em]
    \textbf{Organizing idea:} AI as a \emph{scholarly instrument} \textbf{and} as a \emph{theological subject} ---
    both ``agentic AI \textbf{for} studying theology'' and ``what theology says \textbf{about} agentic AI.''\\[0.3em]
    \textbf{Three technical pillars:} \textbf{Agentic AI $\cdot$ Retrieval-Augmented Generation (RAG) $\cdot$ Evaluation.}
  \end{tcolorbox}
\end{center}

% ---------- 1. description ----------
\section{Course description}
This seminar teaches students to use \textbf{agentic AI} --- systems that plan, use tools, retrieve
sources, and act over multiple steps --- as instruments of theological and religious-studies
research, while critically examining the theological and ethical questions autonomous AI raises.
Students learn to \emph{design a research process}: an agent that retrieves primary sources
(\textbf{RAG}), cross-references them, drafts with citations, critiques its own output, and --- crucially ---
is \textbf{evaluated} so its claims are checked rather than trusted. Each student (or pair) develops,
executes, evaluates, and defends an original project.

The three pillars build on each other: \textbf{agentic AI} is the system, \textbf{RAG} is how it stays
grounded in real texts, and \textbf{evaluation} is how you know whether any of it is trustworthy ---
especially vital where a fabricated scripture reference is a scholarly (and pastoral) failure.

% ---------- 2. objectives ----------
\section{Learning objectives}
By the end of the course a student can:
\begin{enumerate}[leftmargin=1.6em,itemsep=2pt]
  \item \pillar{Agentic AI} Explain what turns an LLM into an agent (tool use, memory, plan$\rightarrow$act$\rightarrow$observe,
    autonomy) and build a multi-step, self-critiquing research agent.
  \item \pillar{RAG} Build a retrieval-augmented pipeline over a theological corpus so the agent cites
    \emph{real} sources, and diagnose where grounding fails.
  \item \pillar{Evaluation} Design and run an evaluation of a RAG/agent system --- retrieval quality,
    groundedness/faithfulness, answer relevance, task success, and cost --- using golden datasets,
    LLM-as-judge, and human review, while knowing each method's limits.
  \item Design a researchable theological question with a matching agentic method \textbf{and eval plan}.
  \item Analyze the theological questions autonomy raises --- personhood, \emph{imago Dei}, agency,
    responsibility, spiritual authority --- and report AI-assisted scholarship \textbf{honestly}.
\end{enumerate}

% ---------- 3. track ----------
\section{Track decision --- Hybrid (low-code default, optional light-code)}
Low-code by default: students build with provided notebook templates and a config-driven builder,
keeping the focus on theology and research design. A parallel \textbf{light-code (Python) path} offers
the same milestones for technically inclined students. Every lab ships with a working scaffold ---
no one writes agents or eval harnesses from scratch.

% ---------- 4. toolkit ----------
\section{Standardized toolkit}
\begin{tabularx}{\textwidth}{@{}>{\bfseries}l Y Y@{}}
\toprule
\normalfont\textbf{Layer} & \textbf{Default (low-code)} & \textbf{Light-code option} \\
\midrule
Model access & Hosted chat model with tool-use (instructor-provided) & Same, via API \\
Agent scaffold & Notebook templates (``research agent,'' ``researcher + fact-checker'') & Minimal Python agent loop \\
Retrieval (RAG) & Prebuilt vector-search over the class corpus & Student-assembled components \\
\textbf{Evaluation} & Provided eval notebook (golden Q\&A set, groundedness + citation checks, LLM-as-judge template) & Same, extensible (e.g. a RAG-eval library) \\
Corpus/sources & Curated digital primary sources (see Week 3) & Same \\
Documentation & Shared trace log + prompt/tool/eval journal & Same \\
\bottomrule
\end{tabularx}

\begin{tcolorbox}[colback=light,colframe=accent,boxrule=0.5pt,arc=2pt,left=8pt,right=8pt,top=5pt,bottom=5pt]
\small\textbf{Guest session recommended} around Week 5--6: ``How real agentic systems are built
\emph{and evaluated}'' from a practitioner.
\end{tcolorbox}

% ---------- 5. session shape ----------
\section{Weekly session shape (100 min)}
$\sim$40 min concept \& discussion $\cdot$ $\sim$40 min hands-on lab $\cdot$ $\sim$20 min project touchpoint
(peer review or mentor check-in). Continuous mentoring, not just in Phase 4.

% ==================== PHASE 1 ====================
\phase{Phase 1 --- Foundations (Weeks 1--3)}

\begin{week}{Week 1}{Framing: ``AI for Theology''}
  \item \textbf{Objective:} Locate the field; articulate a personal research interest.
  \item \textbf{Topics:} Instrument vs. subject; digital humanities, computational text study, AI ethics;
    survey of the AI/religion intersection; preview the three pillars.
  \item \textbf{Hands-on:} Map ``questions I care about''; a deliberately naive chatbot query on a doctrine
    question --- note what feels off (previews \emph{why} we need RAG + evaluation).
  \item \textbf{Deliverable:} One-paragraph interest statement.
  \item \textbf{Readings:} A survey on AI \& religion (e.g. Beth Singler); course intro notes.
\end{week}

\begin{week}{Week 2}{From LLMs to \emph{agents} \pillar{Agentic AI}}
  \item \textbf{Objective:} Build the core mental model of an agent.
  \item \textbf{Topics:} Next-token prediction, embeddings, hallucination; the step to agency --- tool use,
    memory, the plan$\rightarrow$act$\rightarrow$observe loop, autonomy.
  \item \textbf{Hands-on:} Watch an agent decompose a question and call a tool; diagnose where it errs.
  \item \textbf{Readings:} Wolfram, \emph{What Is ChatGPT Doing \dots?} (selections); Anthropic, ``Building Effective
    Agents'' (2024); Weng, ``LLM-Powered Autonomous Agents'' (2023). \emph{Optional:} Yao et al., ``ReAct.''
\end{week}

\begin{week}{Week 3}{Sources and the theological corpus}
  \item \textbf{Objective:} Understand the data landscape and its theological freight.
  \item \textbf{Topics:} Digital primary sources (original-language scripture, patristics, Qur'anic corpora,
    Talmud, translations, hymnody, sermons); copyright, canon, translation politics; ``the data is
    theologically loaded'' --- and why that shapes both retrieval and evaluation.
  \item \textbf{Hands-on:} Locate and load the corpus each student will use all term.
  \item \textbf{Readings:} Claire Clivaz on digital humanities \& biblical studies (selection).
\end{week}

% ==================== PHASE 2 ====================
\phase{Phase 2 --- Technical Core: Agentic AI $\cdot$ RAG $\cdot$ Evaluation (Weeks 4--7)}

\begin{week}{Week 4}{Retrieval-Augmented Generation \pillar{RAG}}
  \item \textbf{Objective:} Make an agent cite \emph{real} texts, not invented ones.
  \item \textbf{Topics:} Embeddings and semantic search; chunking; the retrieve$\rightarrow$augment$\rightarrow$generate loop; RAG as
    the agent's key tool; failure modes (bad chunking, wrong retrieval, ignored context).
  \item \textbf{Hands-on:} Build an ``ask-your-corpus'' RAG agent over a Week-3 text.
  \item \textbf{Readings:} Lewis et al., ``Retrieval-Augmented Generation for Knowledge-Intensive NLP Tasks''
    (2020, the idea, not the math).
\end{week}

\begin{week}{Week 5}{Building a research agent \pillar{Agentic AI}}
  \item \textbf{Objective:} Turn a RAG tool into a multi-step, self-critiquing research agent.
  \item \textbf{Topics:} Agent design patterns --- planning, reflection/self-critique; a two-agent setup
    (``researcher'' + ``critic/fact-checker'') to catch fabricated references; text-analysis techniques
    (semantic search, concept-tracing of \emph{hesed}/\emph{agape}/\emph{grace}) reframed as \textbf{tools the agent wields}.
  \item \textbf{Hands-on:} Extend Week-4's RAG into a planning + reflection agent; add a fact-checker.
  \item \textbf{Readings:} Shinn et al., ``Reflexion'' (2023, concept); Wang et al., ``A Survey on LLM-based
    Autonomous Agents'' (2023, selections). \emph{Guest slot.}
\end{week}

\begin{week}{Week 6}{Evaluation: how do you know it works? \pillar{Evaluation} \newtag}
  \item \textbf{Objective:} Measure whether a RAG/agent system is trustworthy.
  \item \textbf{Topics:}
    \begin{itemize}[leftmargin=1.2em,itemsep=1pt,topsep=1pt]
      \item \textbf{RAG evaluation} --- retrieval precision/recall, \textbf{groundedness/faithfulness} (is the answer
        supported by retrieved text?), answer relevance, and \textbf{citation verification} (does the cited
        verse/passage actually exist and say that?).
      \item \textbf{Agent evaluation} --- task success, trajectory / tool-use correctness, cost \& latency
        efficiency, robustness.
      \item \textbf{Methods} --- golden/reference datasets, \textbf{LLM-as-judge} and its pitfalls (bias, inconsistency),
        human evaluation, and when to use which.
      \item \textbf{Theology-specific angles} --- doctrinal fidelity, translation sensitivity, and \textbf{bias across
        traditions} (whose theology is the ``correct answer''?).
    \end{itemize}
  \item \textbf{Hands-on:} Build a small golden Q\&A set for your corpus; run groundedness + citation-existence
    checks and an LLM-as-judge pass on Week-5's agent; report a scorecard.
  \item \textbf{Readings:} Es et al., ``RAGAS'' (2023); Liu et al., ``G-Eval'' (2023); Zheng et al., ``Judging
    LLM-as-a-Judge / MT-Bench'' (2023, selections).
\end{week}

\begin{week}{Week 7}{Research design \& proposal \quad (midpoint gate)}
  \item \textbf{Objective:} Convert curiosity into a researchable question + agent design + \textbf{eval plan}.
  \item \textbf{Topics:} Scope, method, evidence, what counts as a finding; which agent steps and tools; how
    you will \textbf{evaluate} and verify claims (metrics, golden set, human check); reproducibility and
    documenting agentic-AI use (prompts, tools, traces, eval).
  \item \textbf{Deliverable:} 1--2 page proposal + 5-min pitch. \textbf{Must include an evaluation plan.}
  \item \textbf{Readings:} Research-methods primer (notes); one exemplar agentic-DH study.
\end{week}

% ==================== PHASE 3 ====================
\phase{Phase 3 --- Domain Deep-Dives + Project Launch (Weeks 8--10)}

\begin{week}{Week 8}{AI and hermeneutics}
  \item \textbf{Objective:} Interrogate whether an agent can ``interpret.''
  \item \textbf{Topics:} Multi-step reasoning vs. genuine understanding; authorial intent, tradition, the
    interpretive community; higher stakes when an agent \emph{acts} on interpretation.
  \item \textbf{Hands-on:} Projects formally begin; first evaluated agent runs on the student's question.
  \item \textbf{Readings:} Gadamer, \emph{Truth and Method} (selections); a piece on machine ``understanding.''
\end{week}

\begin{week}{Week 9}{Theology \emph{of} agentic AI}
  \item \textbf{Objective:} Examine personhood and autonomy theologically.
  \item \textbf{Topics:} \emph{Imago Dei}, personhood, consciousness; \textbf{agency and autonomy} --- moral status,
    responsibility, delegation when a system \emph{acts}; AI and idolatry; eschatology.
  \item \textbf{Hands-on:} Mentored project check-in.
  \item \textbf{Readings:} Herzfeld, \emph{The Artifice of Intelligence} (2023) selections; Dorobantu on AI \&
    personhood; Vatican, \emph{Rome Call for AI Ethics} (2020) / \emph{Antiqua et Nova} (2025) selections.
\end{week}

\begin{week}{Week 10}{Ethics, bias, and the limits of autonomy}
  \item \textbf{Objective:} Adopt a critical, human-in-the-loop stance (connects back to evaluation).
  \item \textbf{Topics:} Training-data bias (whose theology dominates?); fabricated citations; agentic-specific
    risks --- compounding errors across steps, over-trust in an autonomous ``researcher,''
    pastoral/spiritual-authority harms; designing human oversight; \textbf{bias-aware evaluation}.
  \item \textbf{Hands-on:} Each project adds an explicit verification/oversight step to its eval plan.
  \item \textbf{Readings:} Bender et al., ``On the Dangers of Stochastic Parrots'' (2021); Crawford, \emph{Atlas of
    AI} (2021, selections).
\end{week}

% ==================== PHASE 4 ====================
\phase{Phase 4 --- Mentored Research Execution (Weeks 11--13)}

\begin{week}{Week 11}{Build sprint (studio)}
  \item \textbf{Objective:} Execute the agentic pipeline; troubleshoot traces.
  \item \textbf{Format:} Studio --- students build, mentors circulate; structured trace debugging.
\end{week}

\begin{week}{Week 12}{Evaluate \& critique \pillar{Evaluation, applied}}
  \item \textbf{Objective:} Run a real evaluation of your \emph{own} project and defend it.
  \item \textbf{Format:} Students apply Week-6 methods to their systems (scorecard: groundedness, citation
    checks, task success, bias spot-checks); peer-review swap + mentor stress-test of claims and AI-reliance.
  \item \textbf{Deliverable:} Preliminary results + \textbf{evaluation report} (early-warning checkpoint).
\end{week}

\begin{week}{Week 13}{Write-up \& communication}
  \item \textbf{Objective:} Turn results into an honest argument.
  \item \textbf{Topics:} Presenting agentic work with integrity --- goals, tools, traces, \textbf{eval results},
    failures, human oversight. Draft final paper + slides.
\end{week}

\phase{Week 14 --- Symposium}
Final presentations (cohort + invited faculty), Q\&A, and a closing reflection: \textbf{what did the
agent reveal, what did it distort, how well did it hold up under evaluation, and what did this
teach us theologically?}
\begin{itemize}[leftmargin=1.4em,topsep=2pt]
  \item \textbf{Deliverable:} Final paper + presentation.
\end{itemize}

% ==================== ASSESSMENT ====================
\section{Assessment}
\begin{tabularx}{\textwidth}{@{}Y r l@{}}
\toprule
\textbf{Component} & \textbf{Weight} & \textbf{Due} \\
\midrule
Participation \& project touchpoints & 15\% & Weekly \\
Interest statement & 5\% & Week 1 \\
Research proposal + pitch (incl. eval plan) & 20\% & Week 7 (gate) \\
Preliminary results + evaluation report & 15\% & Week 12 \\
Final paper & 30\% & Week 14 \\
Final presentation & 15\% & Week 14 \\
\bottomrule
\end{tabularx}

% ==================== RUBRICS ====================
\section{Rubrics}

\subsection*{Research proposal (Week 7)}
\begin{itemize}[leftmargin=1.4em,itemsep=1pt]
  \item \textbf{Question (20\%)} --- theologically substantive, researchable, well-scoped.
  \item \textbf{Agent design (20\%)} --- steps, tools, and RAG grounding fit the question.
  \item \textbf{Evaluation plan (25\%)} --- credible metrics + golden set / human check; awareness of failure
    modes and bias.
  \item \textbf{Feasibility \& sources (20\%)} --- corpus available; scope achievable.
  \item \textbf{Communication (15\%)} --- clear pitch and writing.
\end{itemize}

\subsection*{Final paper (Week 14)}
\begin{itemize}[leftmargin=1.4em,itemsep=1pt]
  \item \textbf{Scholarly contribution (25\%)} --- a real finding or a well-argued negative result.
  \item \textbf{Method \& agent design (20\%)} --- sound, well-documented, reproducible RAG/agent pipeline.
  \item \textbf{Evaluation \& critical stance (25\%)} --- claims measured \emph{and} checked; groundedness/citations
    verified; bias and limits addressed; human-in-the-loop evident.
  \item \textbf{Theological/ethical reflection (15\%)} --- engages the ``AI as subject'' questions (Weeks 8--10).
  \item \textbf{Communication (15\%)} --- structure, clarity, honest reporting of AI use.
\end{itemize}

\begin{tcolorbox}[colback=light,colframe=accent,boxrule=0.5pt,arc=2pt,left=8pt,right=8pt,top=5pt,bottom=5pt]
\small\textbf{Agentic-AI use disclosure (required in every deliverable).}
An appendix in every submission: models/tools used, key prompts, representative \textbf{traces},
\textbf{evaluation results/scorecard}, where the agent failed, and how the student verified results.
Undisclosed or unevaluated AI reliance is penalized.
\end{tcolorbox}

% ==================== NOTES ====================
\section{Notes for the instructor}
\begin{itemize}[leftmargin=1.4em,itemsep=2pt]
  \item \textbf{Readings are starting points from the instructor's knowledge --- verify current editions,
    availability, and access before distributing.} Swap in tradition-appropriate sources
    (Christian / Jewish / Islamic / comparative) to match the cohort.
  \item \textbf{The three pillars are cumulative} --- RAG (W4) grounds the agent (W5), and evaluation (W6)
    judges both; keep labs building on the \emph{same} corpus and agent so students see the through-line.
  \item \textbf{Evaluation is the highest-leverage theology-specific skill here} --- a fabricated citation is a
    scholarly failure, so budget real time for the golden-set + citation-verification lab.
  \item \textbf{Decide access early:} hosted model + tool-use quota; whether the light-code path needs an
    optional extra lab.
  \item \textbf{Corpus curation is the top prep task} --- a clean, rights-clear, well-scoped corpus makes every
    RAG and eval lab work.
\end{itemize}

\end{document}
